% Chapter 2
% 
\chapter{Fundamentação teórica} % Main chapter title
\label{chap:Chapter2} % For referencing the chapter elsewhere

%-------------------------------------------------------------------------------
%---------
%
\section{Fundamentos do sono} 
\label{sec:chap2_sleepFundaments} %For referencing this section elsewhere

O sono apresenta-se como um processo biológico essencial à manutenção da saúde e ao funcionamento diário. Envolve um conjunto organizado de processos fisiológicos relacionados com o funcionamento cognitivo, a regulação emocional e a saúde física, contribuindo para o bem-estar e para o desempenho das atividades quotidianas (Grandner, 2022; Ramar et al., 2021; Roland, 2023).

A caracterização do sono compreende diferentes dimensões, incluindo a duração, a continuidade, a regularidade, o horário e a organização dos seus ciclos. Segundo Ramar et al. (2021), um sono saudável requer duração suficiente, uma qualidade adequada, um horário apropriado e ausência de perturbações do sono. A compreensão destas dimensões permite enquadrar tanto a avaliação dos padrões individuais como o desenvolvimento de intervenções destinadas à sua melhoria.

\subsection{Funções do sono}
\label{sec:chap2_sleepFunctions}

Como foi referido acima, o sono participa na manutenção de funções cognitivas e fisiológicas necessárias ao funcionamento do organismo. No domínio cognitivo, apresenta uma relação importante com a aprendizagem e com a memória. A consolidação da memória corresponde ao processo através do qual a informação adquirida pode tornar-se mais estável após a aprendizagem. Num estudo experimental, Gais et al. (2006) observaram que dormir após a aprendizagem favoreceu a recordação posterior de informação declarativa, apoiando e reforçando a participação do sono na sua retenção.

No domínio fisiológico, o sono relaciona-se com o funcionamento cardiovascular, a regulação hormonal, o metabolismo e a atividade imunitária. Ao longo do episódio de sono, a atividade destes sistemas modifica-se em associação com os diferentes estados de sono e vigília. Por exemplo, durante o sono não REM (\textit{Rapid Eye Movement}, a última fase do sono), a frequência cardíaca e a pressão arterial tendem a diminuir, apresentando um comportamento distinto durante o sono REM (National Heart, Lung, and Blood Institute [NHLBI]).

A importância do sono manifesta-se também no funcionamento durante a vigília. A obtenção de sono suficiente e de qualidade contribui para a atenção, a memória, o bem-estar emocional e a realização das atividades diárias, incluindo as tarefas profissionais e académicas. A segurança constitui igualmente uma dimensão relevante, nomeadamente no que respeita ao risco de acidentes de viação associado ao sono insuficiente (Centers for Disease Control and Prevention [CDC], 2024).

A duração insuficiente do sono encontra-se associada a consequências adversas para a saúde. Segundo Watson, Badr, et al. (2015), dormir regularmente menos de sete horas por noite na idade adulta está associado a hipertensão arterial, diabetes, doença cardiovascular, acidente vascular cerebral e depressão, entre outros problemas de saúde. Os autores referem também associações com alterações do funcionamento imunitário, pior desempenho e maior risco de erros e acidentes. Estas associações reforçam a relevância da duração do sono para a saúde, mas não significam que todas estas consequências ocorram em cada indivíduo, nem que dormir menos de sete horas seja, por si só, equivalente a ter insónia.

Estas relações evidenciam a importância de considerar o sono na promoção da saúde e do bem-estar. A sua avaliação deve abranger o período de repouso e as repercussões no funcionamento diurno, atendendo às diferentes dimensões que caracterizam um sono saudável (Ramar et al., 2021).

\subsection{Arquitetura e fases do sono}
\label{sec:chap2_sleepCycle}

A arquitetura do sono corresponde à organização temporal dos seus estados e estádios ao longo de um episódio de sono. Os estádios são etapas do sono identificadas por características fisiológicas, como a atividade cerebral, os movimentos oculares e o tónus muscular.

Durante o sono, distinguem-se dois estados principais: o sono com movimentos oculares rápidos, designado por REM (Rapid Eye Movement), e o sono não REM, habitualmente identificado como NREM. Este último subdivide-se nos estádios N1, N2 e N3, diferenciados pelas características da atividade cerebral e por outros sinais fisiológicos (NHLBI, 2022; (Memar \& Faradji, 2018; Patel et al., 2025).

O estádio N1 corresponde à transição entre a vigília e o sono, é um sono muito ligeiro. No estádio N2, o sono já encontra-se estabelecido, com padrões específicos de atividade cerebral, enquanto o estádio N3 corresponde ao sono profundo, também denominado sono de ondas lentas. Estes três estados são os estados NREM. O sono REM caracteriza-se por movimentos oculares rápidos, atividade cerebral elevada e uma redução marcada do tónus muscular, sendo frequentemente associado à ocorrência de sonhos (NHLBI, 2022).

Os estados NREM e REM alternam ao longo da noite em ciclos cuja composição varia. Em termos gerais, o sono N3 apresenta maior expressão na primeira parte da noite, enquanto os períodos de sono REM tendem a ser mais prolongados na parte final. Podem também ocorrer despertares breves entre ciclos. Assim, a arquitetura do sono apresenta uma organização dinâmica, com alterações na distribuição dos estádios ao longo do episódio de sono (NHLBI, 2022).

\subsection{Regulação circadiana e homeostática}
\label{sec:chap2_regulacaoCircadiana}

A regulação do sono pode ser enquadrada pelo modelo dos dois processos proposto por Borbély (1982). Este modelo descreve a interação entre um processo homeostático, denominado processo S, e um processo circadiano, denominado processo C. Ambos contribuem para a propensão para dormir e para a organização temporal dos períodos de sono e vigília.

O processo homeostático representa a necessidade de sono que aumenta durante a vigília e diminui durante o sono. Quanto maior o período de vigília, maior tende a ser a pressão para dormir. A atividade de ondas lentas durante o sono NREM constitui uma expressão fisiológica relevante deste processo: apresenta valores mais elevados após privação de sono e diminui ao longo do episódio de sono (Borbély, 1982; Dijk \& Czeisler, 1995).

O processo circadiano está associado a ritmos endógenos com um período próximo de 24 horas. Estes ritmos contribuem para organizar a distribuição temporal do sono, da vigília e de outras funções fisiológicas. A luz constitui um sinal ambiental importante para a sua sincronização. A informação luminosa recebida pelos olhos participa na regulação do relógio central, localizado no núcleo supraquiasmático do hipotálamo, permitindo ajustar os ritmos internos ao ciclo ambiental de luz e escuridão (NHLBI).

A interação entre os processos homeostático e circadiano ajudam a explicar a variação da propensão para dormir ao longo do dia. Dijk e Czeisler (1995) demonstraram contributos distintos e interativos destes mecanismos para a estrutura e consolidação do sono. Deste modo, a interpretação de um padrão de sono deve considerar o histórico de sono e vigília e a sua relação com o momento do ciclo circadiano.

\subsection{Variabilidade individual}
\label{sec:chap2_variabilidade}

As necessidades de sono variam ao longo da vida e entre indivíduos. As recomendações de duração refletem diferenças entre grupos etários: segundo o CDC (2024), são recomendadas 11 a 14 horas diárias para crianças entre um e dois anos e 10 a 13 horas entre os três e os cinco anos, incluindo sestas. Para crianças entre os seis e os doze anos, a recomendação é de 9 a 12 horas, enquanto para adolescentes entre os treze e os dezassete anos é de 8 a 10 horas. Nos adultos entre os dezoito e os sessenta anos, são recomendadas sete ou mais horas por noite. Estes intervalos constituem referências populacionais que devem ser interpretadas considerando as características individuais (Hirshkowitz et al., 2015; Watson, Badr, et al., 2015).

Outra dimensão da variabilidade individual é o cronótipo, relacionado com a expressão temporal dos ritmos circadianos e com a tendência para horários de sono e vigília mais precoces ou mais tardios. Esta variação distribui-se ao longo de um contínuo. Num estudo populacional, Fischer et al. (2017) identificaram diferenças de cronótipo associadas à idade e ao sexo, observando uma tendência para horários mais tardios durante a adolescência e mais precoces nas idades seguintes.

Os horários de sono observados resultam também da interação entre características biológicas e condições de vida, incluindo exigências profissionais, compromissos sociais e exposição à luz. Por conseguinte, o horário efetivamente praticado pode diferir da tendência temporal do indivíduo, tornando necessária a interpretação dos registos no respetivo contexto (Fischer et al., 2017; NHLBI).

No contexto desta dissertação, estes fundamentos orientam a consideração de múltiplas dimensões na caracterização do sono. A duração, a continuidade, a regularidade e o horário, articulados com as características e o contexto do utilizador, constituem elementos a explorar na análise dos padrões individuais e na conceção do sistema de apoio à intervenção.

\section{Insónia e dificuldades de sono} 
\label{sec:chap2_insomnia} %For referencing this section elsewhere

As dificuldades de sono podem manifestar-se através de problemas em adormecer, despertares durante a noite, despertar precoce ou insatisfação com o sono obtido. Estas manifestações podem ocorrer pontualmente ou apresentar um padrão persistente. A sua interpretação exige considerar a frequência, a duração, as circunstâncias em que surgem e as repercussões no funcionamento diário (Riemann et al., 2023).

Atualmente, existem várias perturbações que afetam a qualidade do sono, sendo a insónia, apneia obstrutiva do sono (AOS) e o síndrome das pernas inquietas (SPI) das mais comuns a nível mundial (McArdle et al., 2022). 

\subsection{Conceito e manifestações da insónia}
\label{sec:chap2_insomniaManifest}

A insónia caracteriza-se por dificuldades em iniciar ou manter o sono, ou por despertar antes do horário pretendido com dificuldade em voltar a adormecer. Estas dificuldades ocorrem apesar de existir oportunidade e condições adequadas para dormir e estão associadas a insatisfação com o sono e comprometimento do funcionamento diurno (Riemann et al., 2023).

A caracterização da insónia crónica considera a persistência das dificuldades noturnas e das consequências diurnas, com uma frequência de pelo menos três vezes por semana durante, pelo menos, três meses. Estes elementos integram uma avaliação clínica mais ampla, que considera outras condições capazes de explicar ou acompanhar as queixas (Riemann et al., 2023).

Podem também ocorrer dificuldades de curta duração, por exemplo devido a situações de elevado stress ou alterações de rotinas e dos horários. A existência de algumas noites de sono insatisfatórias não equivale, por si só, a uma perturbação de insónia crónica (NHLBI).

Importa igualmente distinguir insónia de sono insuficiente por falta de oportunidade para dormir. Uma pessoa pode dormir poucas horas devido aos seus compromissos, sem apresentar dificuldade em adormecer quando dispõe de oportunidade adequada. Por isso, a duração do sono, isoladamente, é insuficiente para caracterizar a insónia; devem ser consideradas as queixas, as condições de sono e as repercussões diurnas (Riemann et al., 2023).

\subsection{Prevalência e relevância epidemiológica}
\label{sec:chap2_prevalencia}

A interpretação da prevalência da insónia exige atenção às definições e aos métodos utilizados. Estudos que avaliam sintomas ocasionais, resultados de questionários ou perturbações clinicamente caracterizadas podem produzir estimativas diferentes. Consequentemente, os valores devem ser apresentados em conjunto com a população e os critérios considerados.

Benjafield et al. (2025) estimaram uma prevalência global de insónia clinicamente relevante de 16,2\% na população adulta, correspondente a aproximadamente 852 milhões de pessoas. A análise incluiu 18 estudos, com 262 582 participantes, e extrapolou estimativas para países sem dados nacionais adequados. O estudo identificou uma prevalência superior no sexo feminino. Estes resultados evidenciam a dimensão do problema, devendo ser interpretados à luz da heterogeneidade das definições e da disponibilidade desigual de dados entre países. Também cerca de 25\% das crianças e 18\% de adolescentes sofrem de insónia (Trosman \& Ivanenko, 2021).

No contexto português, Oliveira et al. (2024) analisaram dados de 14 341 adultos provenientes do Inquérito Nacional de Saúde de 2019. Os autores identificaram uma prevalência de 42,9\% de dificuldades de sono auto reportadas, superior nas mulheres. A medida baseou-se numa pergunta sobre dificuldade em adormecer, dormir pouco ou dormir demasiado nas duas semanas anteriores. Este valor caracteriza um conjunto amplo de queixas e não corresponde à prevalência de insónia clinicamente diagnosticada.

\subsection{Fatores envolvidos no desenvolvimento e manutenção da insónia}
\label{sec:chap2_fatoresInsonia}

O desenvolvimento e a persistência da insónia podem ser compreendidos através do modelo dos três fatores, também conhecido como modelo dos 3P. Este enquadramento distingue fatores predisponentes, precipitantes e perpetuantes, permitindo organizar a influência de características individuais, acontecimentos desencadeantes e processos que mantêm as dificuldades ao longo do tempo (Spielman et al., 1987).

Os fatores predisponentes correspondem a características que aumentam a vulnerabilidade às dificuldades de sono. Os fatores precipitantes representam acontecimentos ou alterações que podem desencadear essas dificuldades, como situações de stress. Os fatores perpetuantes incluem comportamentos e processos cognitivos que contribuem para a sua manutenção, mesmo depois de o acontecimento inicial ter diminuído ou desaparecido (Spielman et al., 1987).

Neste enquadramento, estratégias adotadas para compensar noites insatisfatórias podem, em determinadas circunstâncias, contribuir para a persistência do problema. Um exemplo é o aumento excessivo do tempo passado na cama, que pode favorecer períodos prolongados de vigília nesse contexto. A preocupação com o sono e as associações estabelecidas entre a cama e a experiência de permanecer acordado constituem igualmente aspetos relevantes para compreender a manutenção das dificuldades (Spielman et al., 1987).

A identificação destes fatores constitui uma base conceptual importante para a \textbf{Terapia Cognitivo-Comportamental} para a Insónia, cuja intervenção incide particularmente nos processos que perpetuam o problema (Walker et al., 2022).


\subsection{Distinção face a outras dificuldades e perturbações do sono}
\label{sec:chap2_outrasDificuladades}

Algumas manifestações associadas à insónia podem também ocorrer noutras condições. A apneia obstrutiva do sono caracteriza-se por episódios de obstrução das vias aéreas superiores durante o sono, que reduzem ou interrompem o fluxo de ar. Pode acompanhar-se de ressonar, pausas respiratórias e sonolência diurna, comprometendo a continuidade e a qualidade do sono (NHLBI).

A síndrome das pernas inquietas caracteriza-se por uma necessidade intensa de mover as pernas, frequentemente acompanhada por sensações desagradáveis. Os sintomas surgem ou agravam-se em repouso, tendem a ser mais intensos ao final do dia ou durante a noite e são temporariamente aliviados pelo movimento. Estas manifestações podem dificultar o adormecimento e a continuidade do sono (National Institute of Neurological Disorders and Stroke [NINDS]).

As diferenças entre estas condições mostram que uma queixa de sono insatisfatório pode ter diferentes explicações. No contexto desta dissertação, a delimitação da população e das dificuldades abrangidas pela aplicação deverá refletir essa diversidade, articulando o apoio à CBT-I com critérios de avaliação e acompanhamento apropriados.

\subsection{Impacto no funcionamento diário e no bem-estar}
\label{sec:chap2_impacto}

As consequências da insónia ultrapassam o período noturno. Podem incluir fadiga, dificuldades de atenção ou concentração, irritabilidade e redução da energia, afetando o desempenho das atividades quotidianas. A presença destas repercussões constitui uma dimensão importante da caracterização do problema e da avaliação das necessidades da pessoa (NHLBI; Walker et al., 2022).

Na perspetiva da terapia ocupacional, o sono integra o domínio das ocupações, entendidas como as atividades que fazem parte da vida quotidiana. A sua relação com o funcionamento diurno pode ser compreendida considerando as rotinas de auto cuidado, o trabalho ou estudo, o lazer e a participação social. Ho e Siu (2018) salientam a importância de considerar conjuntamente a pessoa, o ambiente e a organização das atividades diárias na compreensão e intervenção sobre o sono.

Esta perspetiva permite enquadrar o sono na distribuição do tempo e das atividades ao longo do dia. Ho e Siu (2018) destacam a organização das rotinas e dos padrões de atividade como componentes relevantes da intervenção, atendendo à relação entre as atividades diurnas e o sono. Deste modo, compreender as dificuldades de sono implica considerar o contexto quotidiano em que ocorrem e as necessidades individuais de participação nas diferentes atividades.

A dimensão do bem-estar encontra-se igualmente refletida no estudo português de Oliveira et al. (2024), no qual as dificuldades de sono foram associadas a menor satisfação com a vida. Por se tratar de um estudo transversal, os resultados demonstram uma associação, sem estabelecer uma relação causal ou a sua direção.

Para esta dissertação, este enquadramento fundamenta uma avaliação que articule as manifestações noturnas com a experiência subjetiva e o funcionamento diurno. Além dos indicadores de sono, será pertinente considerar a fadiga percebida, a interferência das dificuldades de sono nas atividades quotidianas e a adesão às atividades propostas pela aplicação. Esta abordagem poderá contribuir para compreender a utilidade do sistema na perspetiva do utilizador, complementando a análise dos padrões de sono.



\section{Avaliação e monitorização do sono} 
\label{sec:chap2_sleepMonitoring} %For referencing this section elsewhere

A avaliação do sono permite caracterizar as dificuldades apresentadas, compreender os padrões individuais e acompanhar alterações ao longo do tempo. No contexto da insónia, devem ser consideradas as manifestações noturnas, a satisfação com o sono e as repercussões no funcionamento diurno. A avaliação clínica pode ser apoiada por questionários, diários de sono e, quando indicado, métodos instrumentais, como a actigrafia e a polissonografia (Riemann et al., 2023).

Estes métodos fornecem informação de natureza distinta. Os instrumentos de autorrelato permitem conhecer a experiência percebida pela pessoa, enquanto os dispositivos recolhem sinais a partir dos quais podem ser medidos ou estimados parâmetros do sono. A interpretação deve, por isso, atender ao que cada método avalia e às suas limitações.


\subsection{Questionários de avaliação do sono e da insónia}
\label{sec:chap2_questionarioSono}

Os questionários permitem organizar e quantificar diferentes dimensões das dificuldades de sono. Entre os instrumentos utilizados encontram-se o \textit{Pittsburgh Sleep Quality Index }(PSQI), dirigido à qualidade do sono, e o \textit{Insomnia Severity Index} (ISI), centrado na gravidade percebida da insónia.

O PSQI avalia a qualidade e as dificuldades de sono durante o mês anterior. Inclui 19 itens de auto resposta que originam sete componentes: qualidade subjetiva do sono, latência de adormecimento, duração, eficiência habitual, perturbações do sono, utilização de medicação para dormir e disfunção diurna. A pontuação global varia entre 0 e 21, correspondendo os valores mais elevados a maiores dificuldades (Buysse et al., 1989).

Este instrumento proporciona uma caracterização abrangente, mas não constitui uma medida específica da perturbação de insónia nem permite, isoladamente, estabelecer esse diagnóstico (Riemann et al., 2023).

O ISI constitui um instrumento breve desenvolvido para avaliar a gravidade da insónia e acompanhar os resultados do tratamento (Bastien et al., 2001). Integra sete itens relativos às dificuldades em iniciar e manter o sono, ao despertar precoce, à satisfação com o sono, à interferência no funcionamento diário, à perceção do problema por outras pessoas e à preocupação associada. O período habitual de referência corresponde às duas semanas anteriores. Cada item é pontuado entre 0 e 4, resultando numa pontuação global entre 0 e 28, com valores superiores a indicar maior gravidade percebida (Morin et al., 2011).

A escolha entre estes instrumentos depende da dimensão que se pretende avaliar. Para esta dissertação, o PSQI poderá apoiar a caracterização geral da qualidade do sono, enquanto o ISI poderá ser considerado para acompanhar a evolução dos sintomas de insónia. A calendarização das avaliações deverá respeitar os períodos de referência dos instrumentos, distinguindo-os dos registos diários.

\subsection{Diário de sono e indicadores de acompanhamento}
\label{sec:chap2_diarioSono}

O diário de sono consiste no registo prospetivo de informação sobre cada episódio de sono. O \textit{Consensus Sleep Diary}, apresentado por Carney et al. (2012), procura uniformizar os elementos recolhidos e as instruções de preenchimento, favorecendo a comparação dos registos. Inclui horários, estimativas de adormecimento e despertares, bem como uma apreciação subjetiva da qualidade do sono.

A orientação europeia recomenda que a avaliação da insónia seja apoiada por um diário preenchido durante um período de 7 a 14 dias (Riemann et al., 2023). A recolha de várias noites permite observar a variabilidade dos padrões, complementando uma apreciação global.

Entre os indicadores habitualmente considerados encontram-se os seguintes (Carney et al., 2012):
\begin{itemize}
	\item Latência de adormecimento (SOL) - Tempo estimado entre a tentativa de adormecer e o início do sono;
	\item Vigília após o início do sono (WASO) - Tempo de vigília acumulado após o adormecimento, segundo a convenção adotada;
	\item Tempo total de sono (TST) - Duração total estimada do sono;
	\item Tempo na cama (TIB) - Intervalo passado na cama, segundo os horários registados;
	\item Eficiência do sono (SE) - Percentagem do tempo na cama que corresponde a sono;
	\item Qualidade subjetivo do sono - Apreciação da pessoa sobre o sono dessa noite.
\end{itemize}

A eficiência do sono pode ser expressa por:

\[EficiênciaSono (\%) = (\frac{TST}{TIB}) * 100\]

As definições utilizadas no cálculo devem ser explicitadas e mantidas ao longo da avaliação, nomeadamente quanto ao tratamento do despertar final e dos períodos de vigília (Carney et al., 2012).

\subsection{Polissonografia e actigrafia}
\label{sec:chap2_polissonografia}

A polissonografia permite registar simultaneamente diferentes sinais fisiológicos durante o sono. Estes podem incluir atividade cerebral, movimentos oculares, atividade muscular, frequência cardíaca, respiração e saturação de oxigénio. A combinação dos sinais possibilita uma caracterização detalhada do sono e a investigação de determinadas perturbações (National Heart, Lung, and Blood Institute [NHLBI], 2022).

Apesar dessa capacidade, a polissonografia não é necessária na avaliação de rotina de todos os casos de insónia. Pode ser indicada quando existe suspeita de outras perturbações, como alterações respiratórias durante o sono. A caracterização da insónia continua a depender do historial clínico e das queixas apresentadas (Riemann et al., 2023).

A actigrafia utiliza um dispositivo, habitualmente colocado no pulso, para registar o movimento ao longo do tempo e estimar padrões de sono e vigília. Permite acompanhar vários dias no contexto quotidiano. A\textit{ American Academy of Sleep Medicine} apresenta uma recomendação condicional para a sua utilização na estimativa de parâmetros do sono em adultos com insónia, salientando que a monitorização objetiva não é obrigatória para o diagnóstico de rotina (Smith et al., 2018).

Importa distinguir a actigrafia clínica das funcionalidades de monitorização disponibilizadas por dispositivos de consumo. A orientação de Smith et al. (2018) não abrange estes últimos, pelo que as suas recomendações não devem ser automaticamente transferidas para qualquer \textit{smartwatch}.

\subsection{Monitorização através de smartwatches}
\label{sec:chap2_monitoringSmartWatch}

Os \textit{smartwatches} oferecem a possibilidade de acompanhamento do sono em contexto quotidiano. Dependendo do equipamento, os \textit{smartwatches} podem combinar informação de movimento e sinais cardíacos para produzir estimativas de duração, períodos de vigília e estádios do sono. Estes resultados dependem dos sensores e dos algoritmos utilizados pelo fabricante, devendo ser interpretados apenas como estimativas (Chinoy et al., 2021).

Num estudo com 34 adultos jovens saudáveis, Chinoy et al. (2021) compararam sete dispositivos de consumo com a polissonografia. Os equipamentos apresentaram maior capacidade para identificar sono do que vigília, e a classificação dos estádios revelou resultados inconsistentes. O desempenho também variou entre dispositivos e condições de sono. Estes resultados referem-se aos modelos e versões estudados, não permitindo atribuir uma fiabilidade universal aos \textit{smartwatches}.

Para esta dissertação, estas limitações justificam a avaliação das capacidades concretas do equipamento selecionado. Será necessário identificar os dados acessíveis, a sua resolução temporal, as condições de recolha e o momento em que ficam disponíveis para a aplicação. A existência de um resumo de sono após o despertar não demonstra, por si só, que seja possível identificar um despertar noturno no momento em que ocorre.

\subsection{Integração das fontes de informação}
\label{sec:chap2_integration}

Os questionários, o diário e os dispositivos devem ser considerados fontes complementares. A experiência percebida pelo utilizador deve constituir a dimensão central da avaliação da insónia, mesmo quando não coincide integralmente com os registos instrumentais (Riemann et al., 2023).

No contexto desta dissertação, propõe-se que os questionários apoiem a caracterização inicial e a avaliação da evolução, que o diário permita acompanhar a experiência de cada noite e que o \textit{smartwatch} complemente essa informação com estimativas recolhidas automaticamente.

Esta integração deverá preservar a origem de cada registo, distinguindo os dados auto reportados, sinais recolhidos e indicadores calculados. Deverá também considerar a existência de períodos sem dados e eventuais divergências entre as fontes. A ausência de um registo do dispositivo não deverá ser interpretada automaticamente como ausência de sono ou de dificuldades.

Por fim, a avaliação do sistema deverá distinguir a qualidade das estimativas do \textit{smartwatch} da evolução dos sintomas e da utilidade das recomendações. Esta distinção permitirá analisar separadamente a fiabilidade da monitorização e o contributo da intervenção para o utilizador.

\section{Terapia Cognitivo-Comportamental para a Insónia} 
\label{sec:chap2_tcc} %For referencing this section elsewhere

Na intervenção no sono, o papel da Terapia Ocupacional pode estar presente de diversas formas, tais como a adaptação do ambiente, estabelecer novas rotinas antes de dormir ou ainda a utilização de produtos de apoio (Faulkner, 2022; Ho \& Siu, 2018). Para além disso, é importante destacar a Terapia Cognitivo-Comportamental para a Insónia (TCCi) como ferramenta de intervenção nesta área (Faulkner, 2022; Ho \& Siu, 2018). 

A Terapia Cognitivo-Comportamental para a Insónia (TCC-I), designada em inglês por \textit{Cognitive Behavioral Therapy for Insomnia} (CBT-I), é uma intervenção estruturada que combina estratégias cognitivas e comportamentais. A orientação europeia recomenda-a como tratamento de primeira linha para a insónia crónica em adultos, incluindo pessoas com condições clínicas concomitantes (Riemann et al., 2023). A American Academy of Sleep Medicine apresenta igualmente uma recomendação forte a favor da utilização da TCC-I multicomponente nesta população (Edinger et al., 2021).

\subsection{Fundamentos e objetivos da intervenção}
\label{sec:chap2_fundamentoTCC}

A TCC-I incide particularmente nos comportamentos e processos cognitivos que perpetuam a insónia, articulando-se com o modelo dos 3P apresentado anteriormente no subcapítulo ~\ref{sec:chap2_fatoresInsonia}. Este método de intervenção procura atuar sobre comportamentos e fatores mal adaptativos impulsionadores da insónia, através de diferentes fases, como limitar o tempo passado na cama, evitar realizar atividades na cama que não estejam relacionadas com o sono, terapia cognitiva e educação terapêutica sobre higiene do sono. A pessoa participa ativamente na intervenção, através da observação dos seus padrões e da aplicação das estratégias acordadas. Procura de igual forma modificar respostas que contribuem para a manutenção das dificuldades, mesmo quando o acontecimento que inicialmente as desencadeou já desapareceu (Walker et al., 2022).

\subsection{Principais componentes da TCC-I}
\label{sec:chap2_componenteTCC}

A TCC-I integra várias componentes que podem ser articuladas de acordo com o protocolo e as necessidades identificadas. Entre as principais encontram-se a psicoeducação, o controlo de estímulos, a restrição do tempo na cama, a intervenção cognitiva e o relaxamento.

\subsection*{Psicoeducação sobre o sono}

A psicoeducação proporciona informação sobre o funcionamento do sono, as alterações relacionadas com a idade e os mecanismos da sua regulação. A explicação dos processos homeostático e circadiano permite compreender a fundamentação das estratégias propostas, favorecendo uma participação informada no tratamento (Riemann et al., 2023).

\subsection*{Controlo de estímulos}

O controlo de estímulos procura reforçar a associação entre a cama e o sono. Inclui orientações como deitar apenas quando existe sonolência, sair da cama quando não se consegue dormir e reservar a sua utilização unicamente para o sono e a atividade sexual. Estas estratégias integram um plano destinado a reorganizar a relação entre o ambiente de dormir e o comportamento de sono (NHLBI, 2022).

\subsection*{Restrição do tempo na cama}

A restrição do sono, também descrita como restrição do tempo na cama, procura reduzir o desajustamento entre a oportunidade de dormir e o tempo efetivamente dormido. O intervalo inicialmente definido é ajustado ao longo do tratamento, procurando consolidar o sono e, posteriormente, proporcionar uma duração adequada às necessidades individuais. A configuração desta estratégia varia entre protocolos (Riemann et al., 2023).

A sua aplicação corresponde, portanto, a uma intervenção planeada e acompanhada, com critérios de ajuste. A descrição desta componente não deve ser interpretada como uma recomendação geral para reduzir o número de horas de sono.

\subsection*{Intervenção cognitiva}

A intervenção cognitiva procura identificar e avaliar pensamentos e crenças relacionados com o sono que contribuem para a preocupação e a ativação antes de dormir. A reestruturação cognitiva permite examinar a sua adequação e desenvolver interpretações mais realistas, através de um processo de descoberta orientada (Walker et al., 2022).

\subsection*{Relaxamento}

As técnicas de relaxamento procuram reduzir a ativação física e mental associada às dificuldades de sono. Podem integrar o plano de tratamento como exercícios a aprender e praticar, complementando as restantes componentes da intervenção (NHLBI, 2022).

\subsection*{Higiene do sono}

A higiene do sono abrange hábitos e condições ambientais relevantes para dormir. Contudo, a sua utilização isolada não corresponde à TCC-I multicomponente. A orientação da \textit{American Academy of Sleep Medicine} recomenda que a higiene do sono não seja utilizada como única intervenção no tratamento da insónia crónica em adultos (Edinger et al., 2021).

Esta distinção é relevante para a caracterização de aplicações móveis: disponibilizar sugestões gerais sobre hábitos de sono não demonstra, por si só, que uma aplicação implemente um programa estruturado de TCC-I.

\subsection{Organização do tratamento e adesão}

O tratamento desenvolve-se habitualmente ao longo de várias sessões. Walker et al. (2022) descrevem um percurso que inclui uma avaliação inicial, recolha de dados através do diário de sono, introdução das estratégias adequadas, revisão da evolução e prevenção de recaídas. A fase final procura consolidar as aprendizagens e preparar a gestão de futuras dificuldades.

A adesão ao programa constitui uma dimensão importante deste processo. Podem surgir dificuldades no preenchimento do diário ou no cumprimento das orientações acordadas. A sua identificação deve conduzir à exploração das barreiras e à procura colaborativa de soluções, evitando uma abordagem de culpabilização da pessoa (Walker et al., 2022).

Para o projeto desta dissertação, esta perspetiva sugere que o acompanhamento da adesão deverá ultrapassar a contagem de acessos à aplicação. Será pertinente distinguir a consulta dos conteúdos, a compreensão das orientações e a sua aplicação no quotidiano.


\subsection{Evidência de eficácia}

A TCC-I apresenta uma eficácia tão elevada quanto a intervenção farmacológica, uma vez que é considerada o tratamento de primeira linha desta condição. Esta eficácia encontra-se sustentada por investigação experimental e sínteses da literatura. Numa revisão sistemática e meta-análise de 20 estudos, com 1 162 participantes, Trauer et al. (2015) avaliaram intervenções presenciais multicomponentes em adultos com insónia crónica, comparativamente a condições de controlo inativas.

Após o tratamento, os autores observaram reduções médias de aproximadamente 19 minutos na latência de adormecimento e 26 minutos na vigília após o início do sono, assim como uma melhoria de cerca de 10 pontos percentuais na eficiência do sono. O aumento médio do tempo total de sono foi mais reduzido e o respetivo intervalo de confiança incluiu a ausência de efeito. Estes resultados mostram que a melhoria pode manifestar-se sobretudo na continuidade e consolidação do sono, sem um aumento imediato expressivo da sua duração (Trauer et al., 2015).

Contudo, é sempre importante referir que a eficácia da intervenção está condicionada pela colaboração e envolvimento do paciente (Faulkner, 2022; Ho \& Siu, 2018; Walker et al., 2022).


\subsection{Adequação clínica e acompanhamento}

A aplicação das componentes da TCC-I exige atenção às características individuais de cada utilizador e à evolução durante o tratamento. A restrição do tempo na cama pode aumentar temporariamente a sonolência e a fadiga, justificando a avaliação da adequação da estratégia e o acompanhamento dos seus efeitos no funcionamento diurno, incluindo na condução (Riemann et al., 2023).

No contexto da aplicação proposta, estes aspetos fundamentam a definição de critérios de elegibilidade, limites de atuação e situações de encaminhamento. As decisões sobre componentes que alterem a oportunidade de sono deverão seguir um protocolo clinicamente fundamentado e validado por profissionais com competência na área.

\subsection{Implicações para a aplicação proposta}

A TCC-I constitui o enquadramento terapêutico central da aplicação a desenvolver. A personalização deverá apoiar a aplicação de um programa coerente, relacionando os conteúdos e o acompanhamento com as necessidades e a evolução do utilizador.

Neste sentido, podem ser exploradas funcionalidades como o apoio ao preenchimento do diário, a disponibilização progressiva de conteúdos, o acompanhamento de exercícios e a identificação de dificuldades de adesão. Os dados do \textit{smartwatch} poderão complementar a informação auto reportada, respeitando as limitações discutidas na secção anterior.

O momento de entrega deverá também ser considerado em função da estratégia terapêutica. A explicação de um exercício, a preparação para a rotina de deitar e a revisão de dificuldades de adesão podem corresponder a oportunidades distintas de apoio. A adequação dessas oportunidades deverá ser definida e avaliada no projeto.

A eficácia estabelecida para a TCC-I fundamenta a escolha desta abordagem, mas não demonstra antecipadamente o benefício adicional das notificações personalizadas ou da seleção automática do momento de intervenção. Esse contributo constitui uma questão de investigação da presente dissertação, a avaliar através do protótipo e do estudo proposto.


\section{Intervenções digitais e personalização} 
\label{sec:chap2_digitalInterventions} %For referencing this section elsewhere

O avanço tecnológico atual levou a um aumento no uso diário de dispositivos eletrónicos, especialmente durante o período noturno, o que pode prejudicar o sono devido à exposição à luz desses dispositivos (Grotto et al., 2024; Roland, 2023). No entanto, apesar desta tendência negativa, tecnologias como aplicações móveis podem ser uma solução eficaz para melhorar a qualidade do sono, uma vez que facilitam o acesso a informações sobre hábitos saudáveis e oferecem estratégias para problemas de sono (Chang et al., 2015; Grandner, 2022; Grotto et al., 2024; Lancaster et al., 2023; Orzech et al., 2016; Roland, 2023). Deste modo, vários estudos sugerem o potencial do mHealth como forma de intervenção para problemas de sono e de provocar mudanças comportamentais nos seus utilizadores (Arroyo \& Zawadzki, 2022; Grotto et al., 2024; Lancaster et al., 2023).
 